\documentclass[11pt]{letter}
\usepackage[utf8]{inputenc}
\usepackage[T1]{fontenc}
\usepackage{lmodern}
\usepackage{hyperref}

\signature{Qingyong Wang and Lichuan Gu}
\address{School of Artificial Intelligence\\
Anhui Agricultural University\\
Hefei 230036, China\\
Email: wqy@ahau.edu.cn; glc@ahau.edu.cn\\
\vspace{0.3cm}
Anhui Province Key Laboratory of Intelligent Agricultural Technology and Equipment\\
Anhui Agricultural University\\
Hefei 230036, China}

\begin{document}

\begin{letter}{Editorial Office\\
Briefings in Bioinformatics\\
Oxford University Press\\
Great Clarendon Street\\
Oxford OX2 6DP\\
United Kingdom}

\opening{Dear Editor,}

We are pleased to submit our manuscript entitled:

\textbf{``Paired prediction of GPCR--G protein coupling specificity using protein language models and topology-aware feature engineering''}

for consideration as an \textbf{Original Research Article} in \textit{Briefings in Bioinformatics}.

\textbf{Highlights}
\begin{itemize}
  \item Reformulates GPCR--G protein coupling as a paired binary classification over 1,647 (GPCR, G protein family) pairs---incorporating G protein identity as explicit input rather than a class label.
  \item Cross-attention + ESM-2 650M embeddings + topology-aware ICL features achieves AUC 0.862 (cluster-aware CV); G protein family identity alone (4-d one-hot) matches full 1280-d embedding performance, refining the paired-formulation interpretation.
  \item Critical dimension-alignment requirement empirically confirmed: ICL local features must match global embedding dimension---mixing dimensions degrades performance across all architectures.
  \item Thirty-eight standard AlphaFold structural descriptors provide no incremental benefit beyond ESM-2 650M embeddings ($\Delta$AUC $\le$ 0.002), supporting sequence-only pipelines for family-level coupling classification.
  \item Open-source tool \texttt{gpcr-coupling} with pre-trained models and tutorial; all code, data, and results at \url{https://github.com/nblvguohao/gpcr-coupling-leakage-benchmark}.
\end{itemize}

\textbf{Summary}

G protein-coupled receptors (GPCRs) mediate cellular responses by coupling to heterotrimeric G proteins, yet computational prediction of coupling specificity---which determines downstream signaling---remains an open challenge with direct implications for orphan receptor deorphanization and rational drug design. Existing approaches predominantly treat this as single-protein classification, neglecting the paired nature of the interaction.

In this study, we reformulate GPCR--G protein coupling prediction as a pairwise binary classification problem over 1,647 experimentally validated (GPCR, G protein family) pairs spanning 431 GPCRs across four G protein families. Our approach integrates protein language model embeddings (ESM-2, up to 650M parameters), topology-aware intracellular loop features, and a cross-attention architecture with systematic comparison against five baseline methods.

\textbf{Key contributions:}

\begin{enumerate}
    \item \textbf{Systematic benchmarking:} Comprehensive comparison of SVM, MLP, Random Forest, XGBoost, and cross-attention architectures under three evaluation protocols (random CV, cluster-aware CV, leave-one-G-protein-family-out), with classical baselines (AAC+SVM, k-NN) for reference.

    \item \textbf{Methodological findings:} (i) A dimension alignment requirement for local-global protein language model feature fusion---mixing heterogeneous embedding dimensions systematically degrades performance. (ii) Under the 38 AlphaFold structural descriptors examined, ESM-2 650M embeddings capture overlapping information such that explicit structural features contribute no incremental benefit ($\Delta$AUC $\le$ 0.002), supporting sequence-only pipelines for family-level coupling prediction.

    \item \textbf{Rigorous evaluation:} The cross-attention model achieves AUC 0.862 under strict cluster-aware cross-validation. Under this protocol, both cross-attention and SVM produce moderately calibrated probability estimates (Brier scores 0.163 and 0.135, versus constant-predictor baseline 0.190); confidence-stratified accuracy ranges from 83\% to 89\% across thresholds. LOGPSO evaluation reveals a $\sim$0.60 AUC ceiling across \textit{all} tested architectures, highlighting cross-family generalization as a fundamental challenge for current PLM-based approaches.

    \item \textbf{Practical deployment:} An open-source command-line tool (\texttt{gpcr-coupling}) with pre-trained models and a step-by-step tutorial enables both reproduction and application to novel GPCRs, including a case study on orphan receptor deorphanization.

    \item \textbf{Educational value:} The manuscript includes practical guidelines distilled from the systematic evaluation and a tutorial section, consistent with BIB's mission of providing actionable briefings for the bioinformatics community.
\end{enumerate}

\textbf{Why Briefings in Bioinformatics?}

We believe this work aligns closely with \textit{Briefings in Bioinformatics} for several reasons. First, the comprehensive benchmarking framework serves as a reference for future GPCR coupling prediction studies---a briefing function that the journal's readership values. Second, the methodological findings (dimension alignment, AlphaFold redundancy) are generalizable beyond GPCRs and provide practical guidance for researchers working with protein language model embeddings across scales. Third, the accompanying software tool and tutorial lower the barrier to adoption, consistent with BIB's emphasis on methodological accessibility. Fourth, the LOGPSO benchmark proposal and systematic calibration analysis address gaps in the current literature that the BIB community is well-positioned to advance.

\textbf{Manuscript specifications:}

\begin{itemize}
    \item Main text: $\sim$7,000 words, 6 figures, 3 tables
    \item Supplementary: 8 additional tables (including PRAUC values, confidence stratification, hyperparameter configurations), 3 supplementary figures, parameter sensitivity analysis, AlphaFold feature availability analysis
    \item References: 45
    \item All code, data, and pre-trained models publicly available
\end{itemize}

\textbf{Data and Code Availability}

All source code, curated datasets, pre-computed features, and trained model weights are freely available at \url{https://github.com/nblvguohao/gpcr-coupling-leakage-benchmark} under the MIT license.

\textbf{Author Contributions}

G.L., Y.X., and H.L. conceived the study and designed experiments. G.L. implemented models and performed analyses. H.L. contributed to data curation. X.Z. and S.Y. contributed to methodology refinement and result interpretation. L.G. and Q.W. supervised the research. All authors reviewed and approved the manuscript.

\textbf{Conflict of Interest}

The authors declare no competing financial or non-financial interests.

\textbf{Suggested Reviewers}

\begin{enumerate}
    \item David E. Gloriam (University of Copenhagen) --- GPCRdb lead, expertise in GPCR bioinformatics
    \item Bonnie Berger (MIT) --- Protein language models and protein-protein interaction prediction
    \item Arne Elofsson (Stockholm University) --- Protein structure prediction and bioinformatics methods
    \item Johannes S\"oding (Max Planck Institute) --- Sequence-based protein bioinformatics
\end{enumerate}

\closing{Thank you for considering our submission. We look forward to your response.}

\vspace{0.5cm}

Sincerely,

\vspace{0.5cm}

\textbf{Qingyong Wang and Lichuan Gu}

On behalf of all co-authors

\end{letter}

\end{document}
